\documentclass[12pt, oneside]{book}
\usepackage{../mypackages}
\begin{document}
\frontmatter
\title{{\Huge{\textbf{Homological Algebra}}}}
\maketitle

\dominitoc % 初始化minitoc
\pagenumbering{Roman}
\tableofcontents % 主目录


\mainmatter
\pagenumbering{arabic} % 正文编页码字体 





\chapter{}

\section{Preliminaries} % Preliminaries on Topology
\subsection{Topologies}
\begin{definition}
    A topological space $X$ is called \textbf{irreducible} if it cannot be expressed as the union $X=F_1 \cup F_2$ of two proper closed subsets.

    A subset $Y$ of a topological space $X$ is called irreducible if it is irreducible with respect to the induced topology.

    The empty set is not considered to be irreducible.
\end{definition}

\begin{proposition}
    The following conditions on a topological space $X$ are equivalent:
    \begin{enumerate}
        \item
              $X$ is irreducible.
        \item
              Every nonempty open subset of $X$ is irreducible
        \item
              Every nonempty open subset of $X$ is dense in $X$.
        \item
              Any two nonempty open subsets of $X$ intersect.
        \item
              Every nonempty open subset of $X$ is connected.
    \end{enumerate}
\end{proposition}

\begin{definition}
    A topological space $X$ is called \textbf{noetherian} if it satisfies the descending chain condition for closed subsets: for any sequence $Y_1 \supseteq Y_2 \supseteq \ldots$ of closed subsets, there is an integer $r$ such that $Y_r=Y_{r+1}=\ldots$.
\end{definition}

\begin{proposition}
    In a noetherian topological space $X$, every nonempty closed subset $Y$ can be expressed as a finite union $Y=Y_1 \cup \ldots \cup Y_r$ of irreducible closed subsets $Y_i$. If we require that $Y_i \nsupseteq Y_j$ for $i \neq j$, then the $Y_i$ are uniquely determined. They are called the \textbf{irreducible components} of $Y$.
\end{proposition}

\subsection{Algebra}



\section{Affine Varieties} % Affine Varieties
Let $k$ be a fixed algebraically closed field. We define \textbf{affine $n$-space} over $k$, denoted $\mathbf{A}_k^n$ or simply $\mathbf{A}^n$, to be the set of all $n$-tuples of elements of $k$. An element $P \in \mathbf{A}^n$ will be called a point, and if $P=\left(a_1, \ldots, a_n\right)$ with $a_i \in k$, then the $a_i$ will be called the coordinates of $P$.

Let $A=k\left[x_1, \ldots, x_n\right]$ be the polynomial ring in $n$ variables over $k$. We will interpret the elements of $A$ as functions from the affine $n$-space to $k$, by defining $f(P)=f\left(a_1, \ldots, a_n\right)$, where $f \in A$ and $P \in \mathbf{A}^n$. Thus if $f \in A$ is a polynomial, we can talk about the set of zeros of $f$, namely $Z(f)=\left\{P \in \mathbf{A}^n \mid f(P)=0\right\}$. More generally, if $T$ is any subset of $A$, we define the zero set of $T$ to be the common zeros of all the elements of $T$, namely
\begin{equation*}
    Z(T)=\left\{P \in \mathbf{A}^n \mid f(P)=0 \text { for all } f \in T\right\} .
\end{equation*}
Clearly if $\mathfrak{a}$ is the ideal of $A$ generated by $T$, then $Z(T)=Z(\mathfrak{a})$. Furthermore, since $A$ is a noetherian ring, any ideal a has a finite set of generators $f_1, \ldots, f_r$. Thus $Z(T)$ can be expressed as the common zeros of the finite set of polynomials $f_1, \ldots, f_r$.

So for any subset $Y \subseteq \mathbf{A}^n$, let us define the ideal of $Y$ in $A$ by

\begin{equation*}
    I(Y)=\left\{f \in A \mid f(P)_{-}=0 \text { for all } P \in Y\right\}
\end{equation*}
Now we have a function $Z$ which maps subsets of $A$ to algebraic sets, and a function $I$ which maps subsets of $\mathbf{A}^n$ to ideals. Their properties are summarized in the following proposition.


\begin{proposition}
    \begin{enumerate}
        \item
              If $T_1 \subseteq T_2$ are subsets of $A$, then $Z\left(T_1\right) \supseteq Z\left(T_2\right)$.
        \item
              If $Y_1 \subseteq Y_2$ are subsets of $\mathbf{A}^n$, then $I\left(Y_1\right) \supseteq I\left(Y_2\right)$.
        \item $Z(T_1)\cup Z(T_2) = Z(T_1 T_2)$ and $\bigcap Z(T_\alpha)=Z(\bigcup T_\alpha)$
        \item
              For any two subsets $Y_1, Y_2$ of $\mathbf{A}^n$, we have $I\left(Y_1 \cup Y_2\right)=I\left(Y_1\right) \cap I\left(Y_2\right)$.
        \item
              Hilbert's Nullstellensatz.
              For any ideal $\mathfrak{a} \subseteq A, I(Z(\mathfrak{a}))=\sqrt{\mathfrak{a}}$, the radical of $\mathfrak{a}$.
        \item
              For any subset $Y \subseteq \mathbf{A}^n, Z(I(Y))=\bar{Y}$, the closure of $Y$.
    \end{enumerate}
\end{proposition}


\begin{definition}
    A subset $Y$ of $\mathbf{A}^n$ is an \textbf{algebraic set} if there exists a subset $T \subseteq A$ such that $Y=Z(T)$.
\end{definition}


\begin{definition}
    We define the \textbf{Zariski topology} on $\mathbf{A}^n$ by taking the open subsets to be the complements of the algebraic sets.
\end{definition}

\begin{definition}
    An affine algebraic variety (or simply \textbf{affine variety}) is an irreducible closed subset of $\mathbf{A}^n$.
    An open subset of an affine variety is a quasi-affine variety.

\end{definition}


\begin{corollary}
    There is a one-to-one inclusion-reversing correspondence between algebraic sets in $\mathbf{A}^n$ and radical ideals (i.e., ideals which are equal to their own radical) in $A$, given by $Y \mapsto I(Y)$ and $\mathfrak{a} \mapsto Z(\mathfrak{a})$.

    Furthermore, an algebraic set is irreducible if and only if its ideal is a prime ideal.
\end{corollary}




\begin{definition}
    If $Y \subseteq \mathbf{A}^n$ is an affine algebraic set, we define the \textbf{affine coordinate ring} $A(Y):=A/I(Y)$ of $Y$.
    \begin{remark}
        If $Y$ is an affine variety, then $A(Y)$ is an integral domain. Furthermore, $A(Y)$ is a finitely generated $k$-algebra.

        Conversely, any finitely generated $k$-algebra $B$ which is a domain is the affine coordinate ring of some affine variety.
    \end{remark}
\end{definition}















\end{document}